\scanpage{191}
\mspage{498}
% position: upper continuation
\mbox{u n e}\quad\mbox{v i s i o n}, et encore moins une vision du monde. Cela vous fait une
belle jambe d'être ``convaincu'' mordicus d'une vision qui reste étrangère, incom-
prise. Pour tout dire, cela n'a pas même de sens - ou plus exactement, le vrai
sens de sa ``conviction'' n'est pas plus compris par l'intéressé, que cette vision
qu'il fait mine d'incorporer à son lourd bagage culturel.

% position: upper middle paragraph
Quand la vision est comprise et assimilée, la question même d'une ``preuve''
apparaît étrangement saugrenue - un peu comme de prouver que le ciel est bleu
quand on voit bien qu'il est bleu, ou que le parfum d'une fleur qu'on aime est
bon...
\margin{\figure{\begin{tikzpicture}\draw[gray] (0,0) .. controls (-0.5,-0.4) and (-0.8,-1.1) .. (-0.8,-1.7);\end{tikzpicture}}}
% position: lower main paragraph
($^{112}$) (17 octobre) Mes premières réflexions sur le double aspect ``féminin'' et
``masculin'' sont issues d'une réflexion sur moi-même. C'était vers les débuts
1979, à un moment où j'ignorais encore les mots chinois ``yin'' et ``yang'', et
l'existence d'une sorte de ``philosophie'' subtile du jeu incessant du yin et
du yang, dans la tradition culturelle chinoise. J'ai appris la chose vers la
fin de la même année je crois, par ma fille et surtout par mon gendre Ahmed, qui
commençait alors à s'intéresser à la médecine chinoise, sur laquelle il a accro-
ché fortement dans les années suivantes. La plupart de ce qu'il me disait re-
coupait et confirmait la vision à laquelle j'étais parvenu, chose qui n'avait
rien pour me surprendre. Si surprise il y avait, c'était plutôt dans les quel-
ques cas de ``couples'' où le rôle yin-yang ``naturel'' me semblait renversé, dans
la tradition chinoise. Mon réflexe (fortement ``yang'' en l'occurence!) avait été
une conviction à fleur de peau que ce ``renversement'' devait être dû à une dé-
formation culturelle, sans d'ailleurs aller y regarder de trop près (*) - \unsure{c'étai|c'était}
à un moment où mes gammes passées sur le féminin-masculin me paraissaient bien
lointaines, alors que j'étais engagé dans une méditation autrement plus \unsure{personel|personelle}
% position: footnote rule and footnote
\par\noindent\rule{0.24\linewidth}{0.4pt}\par
(*) Cette réaction d'assurance péremptoire, vis-à-vis d'une tradition millénai-
re qui aurait pu m'inciter à plus de prudence, est celle-là même qui, enfant,
m'a fait récuser la formule (bien compliquée ma foi!) $\inserted{\pi}=3,14\ldots$ enseignée par
les livres, en faveur de $\inserted{\pi}=3$ dont je m'étais convaincu par mes propres moyens.
(Voir la note ``La quadrature du cercle'', n° 69.) Il est vrai que pour cette his-
toire du yin et du yang, j'avais eu ample occasion de me rendre compte à quel
point l'appréhension de la nature du ``féminin'' et du ``masculin'', et de leurs
interrelations, est faussée par des distortions culturelles invétérées, d'une
force considérable. Je ne me rendais pas compte encore, par contre, à quel point
aussi une appréhension précise et délicate de ces relations était chose \unsure{essentie|essentiel}le
dans la pratique de certains arts traditionnels chinois, et poussé à un degré
de grande finesse.
% position: lower left strokes and number
\margin{$-$}
\figure{\begin{tikzpicture}\draw[gray!50] (0,0)--(1.8,0.5);\end{tikzpicture}}
\margin{190}
